\scanpage{104}

un fibré vectoriel de rang 3 $V$ (le fibré des translations de $E$)

\[
0\longrightarrow V\mathbin{\text{---}}\unsure{\mathscr E|\mathcal E}\mathbin{\text{---}}\mathcal O_S\mathbin{\text{---}}1
\]

et ``une'' forme quadratique $Q$ \inserted{lisse} sur $\mathscr E$ définie loc$^{\mathrm t}$ mod \unsure{mult$^{\mathrm n}$|mult.} par une unité [c'est ce que signifie que $\Sigma$ est défini par une fonction polynomiale de degré $\leq2$, et qu'il est \unsure{project$^{\mathrm t}$|projectiv$^{\mathrm t}$} lisse\ldots]

La condition que $\Sigma$ ne soit pas un paraboloïde signifie que $q\mid V$ est \deleted{non dégénérée} lisse. Si les car. résiduelles sont $\ne2$, cela signifie \inserted{aussi} que $\mathscr E=V+V^\perp$, donc $E$ se trouve muni d'une origine $o$ \inserted{permettant de l'identifier à $V$} (le polaire de $H=P(V)$ p.r. à \inserted{la quadrique} $Q$) et $\Sigma$ est stable pour la sym. p.r. à cette origine -- $o$ s'appelle le \uline{centre} de la quadrique affine. Ce centre $o$ \inserted{(section de $E$ sur $S$)} ne rencontre pas $\Sigma$, donc la fonction quadratique peut être normalisée \inserted{loc.} par la condition que sa valeur sur $o$ soit 1 ; ce qui la détermine de

% archive sheet number: 103
